% Source: https://artofproblemsolving.com/community/c3h1936640_solve_the_equation
% Stable identity: aops-topic-1936640-q01
The square roots require $8x+1\ge0$ and $46-10x\ge0$, so $-\dfrac{1}{8}\le x\le\dfrac{23}{5}$.

Trying $x=1$: the left side is $\sqrt9+\sqrt{36}=9$ and the right side is $-1+5+4+1=9$, so $x=1$ is a solution. We show that it is the only one.

Subtract $9$ from both sides. On the left, rationalising each difference,
\[
\sqrt{8x+1}-3=\dfrac{8(x-1)}{\sqrt{8x+1}+3},\qquad \sqrt{46-10x}-6=\dfrac{-10(x-1)}{\sqrt{46-10x}+6}.
\]
On the right, $-x^3+5x^2+4x-8=(x-1)(8+4x-x^2)$. So the equation is equivalent to
\[
(x-1)\left[\dfrac{8}{\sqrt{8x+1}+3}-\dfrac{10}{\sqrt{46-10x}+6}-\left(8+4x-x^2\right)\right]=0 .
\]
On the domain the first fraction is at most $\dfrac{8}{3}$ and the second fraction is positive. Also $8+4x-x^2=12-(x-2)^2$, and on the domain $|x-2|\le\dfrac{13}{5}$ (the point of the domain farthest from $2$ is $x=\dfrac{23}{5}$), so
\[
8+4x-x^2\ge12-\dfrac{169}{25}=\dfrac{131}{25}>\dfrac{8}{3} .
\]
Hence the bracket is negative for every $x$ in the domain, and the only solution is $x=\boxed{1}$.
